% =====================================================================
% 07-OperatorTheorems.tex
% Geometric Hallucination Taxonomy — Mathematical Spine
% Borealis S. Hedling — 2026
% =====================================================================

\section*{07. Operator Theorems}

This module formalizes the core theorems linking geometric instability
to hallucination behavior in the Geometric Hallucination Taxonomy (GHT).
Each theorem establishes a precise implication between operator-level
instability and latent trajectory divergence.

All results rely on the foundations (00) and operator modules (01–05).

% ---------------------------------------------------------------------
\subsection*{7.1 Preliminaries}

Let $x(t)$ be a latent trajectory on a smooth manifold $M$.

Let the operator outputs be:
\[
K_t, \quad
\Delta_\lambda(t), \quad
H_{\gamma_t}, \quad
\tau_t, \quad
F_t.
\]

Let the unified stability region be:
\[
\Omega
= \Omega_{\kappa}
\cap \Omega_{\sigma}
\cap \Omega_{h}
\cap \Omega_{\tau}
\cap \Omega_{f}.
\]

A trajectory is stable if $x(t) \in \Omega$ for all $t$.

% ---------------------------------------------------------------------
\subsection*{7.2 Curvature Instability Implies Divergence}

\begin{theorem}[Curvature Divergence]
If there exists $t_0$ such that
\[
|K_{t_0}| > \kappa_{\mathrm{thr}},
\]
then there exists $\varepsilon > 0$ such that geodesics emanating from
$x(t_0)$ diverge at rate at least $\varepsilon$.

Formally, for sufficiently small $s$,
\[
d(\gamma_1(s), \gamma_2(s)) \ge \varepsilon s,
\]
where $\gamma_1$ and $\gamma_2$ are geodesics in the unstable plane.
\end{theorem}

\begin{proof}[Sketch]
Sectional curvature controls geodesic deviation via the Jacobi equation.
Exceeding the threshold forces exponential or hyperbolic divergence.
\end{proof}

% ---------------------------------------------------------------------
\subsection*{7.3 Spectral Instability Implies Local Expansion or Collapse}

\begin{theorem}[Spectral Cascade]
If
\[
|\Delta_\lambda(t_0)| > \sigma_{\mathrm{thr}},
\]
then the local linearization of the update rule $\Phi$ at $x(t_0)$
exhibits either expansion or collapse.

Formally,
\[
\|A_{t_0+1} v\| \ge c \|A_{t_0} v\|
\]
for some $v$ and constant $c > 1$, or the analogous collapse inequality.
\end{theorem}

\begin{proof}[Sketch]
Eigenvalue drift controls the spectral radius of $A_t$. Large drift
forces the radius to change by a multiplicative factor.
\end{proof}

% ---------------------------------------------------------------------
\subsection*{7.4 Holonomy Drift Implies Rotational Instability}

\begin{theorem}[Holonomy Drift]
If
\[
\|H_{\gamma_{t_0}}\| > h_{\mathrm{thr}},
\]
then parallel transport around $\gamma_{t_0}$ induces rotational or
shearing instability in $T_{x(t_0)}M$.

Formally, there exists $v$ such that
\[
\|H_{\gamma_{t_0}} v\| \ge \varepsilon
\]
for some $\varepsilon > 0$.
\end{theorem}

\begin{proof}[Sketch]
Holonomy measures failure of parallel transport to preserve direction.
Exceeding the threshold forces nontrivial rotation or shear.
\end{proof}

% ---------------------------------------------------------------------
\subsection*{7.5 Metric Rank Loss Implies Topological Collapse}

\begin{theorem}[Rank Collapse]
If
\[
\tau_{t_0} < n,
\]
then $x(t_0)$ lies on the collapse surface
\[
\Sigma = \{ p : \det(g_p) = 0 \}.
\]

Furthermore, distances and angles become ill-defined in a neighborhood of
$x(t_0)$.
\end{theorem}

\begin{proof}[Sketch]
Rank loss implies determinant zero. Degenerate metrics fail to define
inner products, collapsing local geometry.
\end{proof}

% ---------------------------------------------------------------------
\subsection*{7.6 Fusion Dominance Implies Nonlinear Spike}

\begin{theorem}[Fusion Spike]
If
\[
\|F_{t_0}\|_g > f_{\mathrm{thr}},
\]
and nonlinear dominance holds:
\[
\|\gamma \sigma(x(t_0), y(t_0))\|_g
\gg
\|\alpha x(t_0) + \beta y(t_0)\|_g,
\]
then the fused representation exhibits a nonlinear spike.

Formally,
\[
\|F_{t_0}\|_g \ge c
\]
for some $c > f_{\mathrm{thr}}$.
\end{theorem}

\begin{proof}[Sketch]
Dominance forces the nonlinear term to control the output norm.
\end{proof}

% ---------------------------------------------------------------------
\subsection*{7.7 Unified Instability Implies Hallucination}

\begin{theorem}[Unified Hallucination Theorem]
If a latent trajectory satisfies
\[
\exists t_0 : x(t_0) \notin \Omega,
\]
and instability propagates through the interaction map $\mathcal{I}$,
then the trajectory exhibits hallucination behavior.

Formally,
\[
\{ t : x(t) \notin \Omega \} \text{ is infinite}.
\]
\end{theorem}

\begin{proof}[Sketch]
Instability in any operator propagates through $\mathcal{I}$, producing
persistent instability. This matches the definition of hallucination.
\end{proof}

% ---------------------------------------------------------------------
\subsection*{7.8 Operator Interaction Theorem}

\begin{theorem}[Operator Interaction]
Instability in one operator may induce instability in another. For
example:
\[
|K_t| > \kappa_{\mathrm{thr}}
\;\Rightarrow\;
|\Delta_\lambda(t+1)| > \sigma_{\mathrm{thr}},
\]
under appropriate geometric conditions.

Similar implications hold for other operator pairs.
\end{theorem}

\begin{proof}[Sketch]
Curvature blowout distorts local linearization, affecting eigenvalues.
Other interactions follow from geometric coupling.
\end{proof}

% ---------------------------------------------------------------------
\subsection*{7.9 Global Instability Implies Divergence}

\begin{theorem}[Global Divergence]
If instability persists,
\[
\{ t : x(t) \notin \Omega \} \text{ is infinite},
\]
then the latent trajectory diverges from any stable region.

Formally,
\[
\limsup_{t \to \infty} d(x(t), \Omega) > 0.
\]
\end{theorem}

\begin{proof}[Sketch]
Persistent instability prevents re-entry into $\Omega$.
\end{proof}

% =====================================================================
% End of 07-OperatorTheorems.tex
% =====================================================================
